\makeabstract
    {
    Optimizing Synthesis of InP/MnZnS Quantum Dots
    }
    {
    Melody Guo\textsuperscript{1}, Ji Hyun Kim\textsuperscript{1}, Jacob Olshansky\textsuperscript{1}
    }
    {
    \textsuperscript{1} Department of Chemistry, Amherst College, Amherst, MA
    }
    {
    This study aimed to synthesize and characterize manganese (Mn)-doped core/shell indium phosphide/zinc sulfide (InP/ZnS) quantum dots (QDs) with blue and green emission. Mn doping introduces additional optical properties to semiconductor QDs, making these materials of interest for applications in optoelectronics. A major challenge encountered during synthesis was the formation of phosphate byproducts, which interfered with the characterization of the desired QDs. We focused on modifying the synthesis conditions to improve QD formation while minimizing the amount of phosphates.
    }
    {
    InP/ZnS QDs were synthesized using a hot-injection method, with modifications made to reaction conditions between syntheses to optimize QD formation. Distinct samples were synthesized in the green-emitting series. Absorption and emission spectroscopy were used to characterize the QDs and determine their emission behavior, while transmission electron microscopy (TEM) was used to examine the morphology and size of QDs. Changes made between syntheses were assessed based on their effects on product formation and characterization. 
    }
    {
    Sample MG02 served as a control sample for comparing subsequent modifications and showed significant phosphate formation. Repeating the synthesis with a degassed petrolatum jelly (PJ):octadecene (ODE) mixture and increased stirring speed in sample MG03 did not substantially resolve phosphate formation. Changing the inert atmosphere from house nitrogen to Grade 5 argon in sample MG04 resulted in a decrease in the amount of phosphate observed. Further modifications in sample JHK22, based on previously successful synthesis protocols, including consolidation of the precursor injections, use of a fresh tris(diethylamino)phosphine source, and stricter control of heating times, resulted in a further decrease in phosphate formation, although phosphate byproducts remained present. 
    }
    {
    Synthesis modifications for the InP/ZnS QDs procedure demonstrated that reaction conditions, including inert atmosphere, precursor amount, injection strategy, and heating control, can influence phosphate byproduct formation. The decrease in phosphate observed following the use of Grade 5 argon and subsequent protocol modifications suggests that these changes improved the synthesis conditions, although further optimization is necessary to eliminate phosphate formation and obtain consistently characterizable QDs. Uv-Vis/photoluminescence spectroscopy and TEM characterization provided methods for identifying conditions of successful Mn-doped and undoped InP/ZnS QDs.
    }

    \newpage
    